\documentclass[12pt,a4paper]{article}
\usepackage{hyphenat}
% ─── Paquetes ────────────────────────────────────────────────────────────────
\usepackage[utf8]{inputenc}
\usepackage[T1]{fontenc}
\usepackage[spanish]{babel}
\usepackage{lmodern}
\usepackage{geometry}
\usepackage{graphicx}
\usepackage{hyperref}
\usepackage{xcolor}
\usepackage{booktabs}
\usepackage{array}
\usepackage{longtable}
\usepackage{enumitem}
\usepackage{fancyhdr}
\usepackage{titlesec}
\usepackage{tocloft}
\usepackage{parskip}
\usepackage{microtype}
\usepackage{amsmath}
\usepackage{caption}
\usepackage{float}
\usepackage{listings}
\usepackage{pgfgantt}
\usepackage{tikz}
\usetikzlibrary{shapes.geometric, arrows.meta, positioning, fit, backgrounds}

% ─── Geometría ───────────────────────────────────────────────────────────────
\geometry{
  top=2.5cm, bottom=2.5cm,
  left=3cm,  right=2.5cm
}

% ─── Colores ─────────────────────────────────────────────────────────────────
\definecolor{primary}{HTML}{1A3557}
\definecolor{accent}{HTML}{E84545}
\definecolor{lightgray}{HTML}{F5F5F5}
\definecolor{codebg}{HTML}{F0F4F8}
\definecolor{codetext}{HTML}{1A1A2E}

% ─── Hiperlinks ──────────────────────────────────────────────────────────────
\hypersetup{
  colorlinks=true,
  linkcolor=primary,
  urlcolor=accent,
  citecolor=primary,
  pdftitle={ResQ-Net: Sistema de Gestión Humanitaria Distribuida},
  pdfauthor={Grupo 10 --- Introducción a la Ingeniería 2026-I}
}

% ─── Cabecera y pie de página ────────────────────────────────────────────────
\pagestyle{fancy}
\fancyhf{}
\fancyhead[L]{\small\color{primary}\textbf{ResQ-Net} --- Ayuda Humanitaria Distribuida}
\fancyhead[R]{\small\color{primary}Introducción a la Ingeniería · 2026-I}
\fancyfoot[C]{\thepage}
\renewcommand{\headrulewidth}{0.4pt}

% ─── Títulos de sección ──────────────────────────────────────────────────────
\titleformat{\section}{\Large\bfseries\color{primary}}{}{0em}{\thesection\quad}[\titlerule]
\titleformat{\subsection}{\large\bfseries\color{primary!80}}{}{0em}{\thesubsection\quad}
\titleformat{\subsubsection}{\normalsize\bfseries}{}{0em}{\thesubsubsection\quad}

% ─── Código ──────────────────────────────────────────────────────────────────
\lstset{
  backgroundcolor=\color{codebg},
  basicstyle=\ttfamily\small,
  keywordstyle=\color{primary}\bfseries,
  commentstyle=\color{gray},
  stringstyle=\color{accent},
  breaklines=true,
  frame=single,
  framerule=0.4pt,
  rulecolor=\color{primary!30}
}

% ─── Documento ───────────────────────────────────────────────────────────────
\begin{document}

% ╔══════════════════════════════════════════════════════════════════════════╗
% ║                          PORTADA                                         ║
% ╚══════════════════════════════════════════════════════════════════════════╝
\begin{titlepage}
  \centering
  \vspace*{1cm}

  {\color{primary}\rule{\linewidth}{2pt}}
  \vspace{0.5cm}

  {\LARGE\bfseries\color{primary}
    DEPARTAMENTO DE INGENIERÍA DE SISTEMAS\\[0.3cm]
    Introducción a la Ingeniería · 2026-I}

  \vspace{0.5cm}
  {\color{accent}\rule{\linewidth}{1pt}}
  \vspace{1.5cm}

  {\Huge\bfseries\color{primary} ResQ-Net}\\[0.4cm]
  {\Large Sistema de Gestión de Recursos Humanitarios\\
  Descentralizado en Tiempo Real}

  \vspace{1cm}
  {\large\color{gray} Caso 10 --- Ayuda Humanitaria}

  \vspace{2cm}
  \begin{tabular}{ll}
    \toprule
    \textbf{Integrante} & \textbf{Rol principal} \\
    \midrule
    Jose Daniel Castro Villadiego & Arquitectura de sistemas y backend \\
    Santiago Andres García Acevedo & Desarrollo frontend / MVP \\
    Santiago de Jesús Trejos Zuluaga & Algoritmos de triaje y lógica \\
    Yeiner Jesús Hernández bustos & Documentación técnica y ética \\
    Yordi Romario Jiménez López & Integración y pruebas \\
    \bottomrule
  \end{tabular}

  \vspace{1.5cm}
  \begin{tabular}{ll}
    \textbf{Docente:} & Miguel Alberto Caro Álvarez \\
    \textbf{Repositorio:} & \href{https://github.com/yei18/ResQ-Net.git}{\texttt{https://github.com/yei18/ResQ-Net.git}} \\
    \textbf{Fecha de entrega:} & 27 de marzo de 2026 \\
  \end{tabular}

  \vfill
  {\color{primary}\rule{\linewidth}{2pt}}

  \vspace{0.4cm}
  {\small\itshape ``La tecnología es inútil si no sirve a la humanidad.''}
\end{titlepage}

% ─── Tabla de contenidos ─────────────────────────────────────────────────────
\tableofcontents
\newpage

% ╔══════════════════════════════════════════════════════════════════════════╗
% ║  1. DESCRIPCIÓN DEL PROBLEMA                                             ║
% ╚══════════════════════════════════════════════════════════════════════════╝
\section{Descripción del Problema}

\subsection{Contexto del caso}

Cuando un desastre natural golpea una zona densamente poblada, los primeros
sesenta minutos son los más críticos. Los equipos de rescate llegan sin mapas
actualizados, sin saber cuántas personas necesitan medicamentos urgentes, sin
forma de evitar que dos brigadas envíen kits al mismo punto mientras otra zona
queda sin cobertura. La infraestructura de telecomunicaciones colapsa. Las
bases de datos centrales dejan de responder. El caos operativo mata a personas
que podrían haberse salvado.

Según datos de la \textbf{OCHA} (Oficina para la Coordinación de Asuntos
Humanitarios de la ONU), entre 2000 y 2023 los desastres naturales afectaron
a más de 4.200 millones de personas y causaron pérdidas económicas superiores
a 3,6 billones de dólares \cite{ocha2023}. En Latinoamérica, Colombia ocupa
el quinto lugar en número de eventos por año \cite{ungrd2022}, con terremotos
en el eje cafetero, inundaciones en la Costa Caribe e incendios forestales en
la Amazonía.

El problema central no es la falta de recursos: es la incapacidad de
\textbf{coordinarlos en tiempo real cuando la infraestructura falla}.

\subsection{Magnitud del problema}

\begin{itemize}[leftmargin=1.5em]
  \item En el terremoto de Haiti (2010), la duplicación de entregas de víveres
        llegó al \textbf{34\,\%} del total de suministros movilizados
        \cite{schuller2016}.
  \item En las inundaciones del Chocó (2019), brigadas de rescate tardaron
        hasta \textbf{6 horas} en sincronizar su posición por falta de
        comunicación directa \cite{ungrd2022}.
  \item El \textbf{80\,\%} de las muertes evitables en catástrofes ocurre en
        las primeras 4 horas, período en el que la coordinación es más difícil
        \cite{who2022}.
\end{itemize}

\subsection{Población afectada}

\begin{itemize}[leftmargin=1.5em]
  \item \textbf{Directa:} brigadas de rescate, coordinadores de operaciones de
        emergencia, pilotos de drones de entrega.
  \item \textbf{Indirecta:} toda persona atrapada o afectada en la zona del
        desastre, comunidades aledañas que requieren evacuación, voluntarios y
        ONGs que llegan sin visibilidad del terreno.
\end{itemize}

\subsection{Justificación de la solución}

Un sistema distribuido que funcione \textbf{sin conexión a internet central}
---usando redes Mesh, computación en el borde (\textit{fog computing}) y
comunicación por radio cuando sea necesario--- permite que los equipos en
campo compartan su estado, las necesidades detectadas y la ubicación de
suministros, aun cuando la infraestructura convencional ha colapsado. Agregar
un motor de triaje algorítmico que priorice la asignación de recursos escasos
reduce la dependencia de decisiones ad hoc bajo presión extrema.

% ╔══════════════════════════════════════════════════════════════════════════╗
% ║  2. CÓMO SE LLEGÓ A LA SOLUCIÓN                                          ║
% ╚══════════════════════════════════════════════════════════════════════════╝
\section{Cómo se llegó a la Solución}

\subsection{Proceso de ideación}

El equipo aplicó un proceso de cuatro fases inspirado en \textbf{Design
Thinking} \cite{plattner2010}:

\begin{enumerate}[leftmargin=1.5em]
  \item \textbf{Empatizar:} revisamos reportes post-desastre de la Cruz Roja,
        la UNGRD y testimonios de voluntarios colombianos que participaron en
        las inundaciones de Mocoa (2017).
  \item \textbf{Definir:} el problema central es la \textit{falta de
        visibilidad compartida} entre equipos que operan en campo sin
        conectividad estable.
  \item \textbf{Idear:} propusimos tres alternativas (véase siguiente sección).
  \item \textbf{Prototipar y testear:} construimos el MVP mínimo para validar
        la lógica de sincronización y el algoritmo de triaje.
\end{enumerate}

\subsection{Alternativas consideradas y descartadas}

\begin{longtable}{>{\bfseries}p{3.5cm} p{5cm} p{5cm}}
  \toprule
  Alternativa & Ventaja & Razón de descarte \\
  \midrule
  \endfirsthead
  \toprule
  Alternativa & Ventaja & Razón de descarte \\
  \midrule
  \endhead
  App centralizada en la nube &
    Interfaz conocida, fácil de desplegar &
    Requiere internet: inutilizable cuando la infraestructura colapsa \\[4pt]
  Radio UHF / VHF analógica &
    No depende de ninguna red digital &
    Sin capacidad de compartir datos estructurados; requiere operadores
    especializados y coordinación manual \\[4pt]
  \textbf{Red Mesh + fog computing} (elegida) &
    Funciona sin internet central; los nodos se sincronizan entre sí;
    escalable de forma orgánica &
    Mayor complejidad inicial, pero cubre los requisitos críticos del
    problema \\
  \bottomrule
\end{longtable}

\subsection{Criterio de selección}

La solución elegida satisface tres requisitos no negociables identificados
durante la fase de definición: (1) operar sin infraestructura centralizada,
(2) compartir estado en tiempo cuasi-real entre nodos dispersos y (3) tomar
decisiones de asignación de recursos con criterios explícitos y auditables.

% ╔══════════════════════════════════════════════════════════════════════════╗
% ║  3. ARQUITECTURA DEL MVP                                                 ║
% ╚══════════════════════════════════════════════════════════════════════════╝
\section{Arquitectura del MVP}

\subsection{Visión general}

ResQ-Net se compone de tres capas que pueden operar de forma autónoma o
coordinada:

\begin{enumerate}[leftmargin=1.5em]
  \item \textbf{Capa de campo:} dispositivos móviles (tablets, smartphones) en
        manos de las brigadas.
  \item \textbf{Capa de niebla (fog):} nodos locales ---Raspberry Pi o
        laptops--- que actúan como coordinadores de sector.
  \item \textbf{Capa de nube (opcional):} sincronización diferida con un
        servidor central cuando haya conectividad disponible.
\end{enumerate}

\subsection{Diagrama de arquitectura}

\begin{figure}[H]
\centering
\begin{tikzpicture}[
  node distance=1.2cm and 1.8cm,
  every node/.style={font=\small},
  box/.style={rectangle, rounded corners=4pt, draw=primary, thick,
              fill=primary!8, minimum width=2.8cm, minimum height=1cm,
              text centered, text width=2.6cm},
  cloud/.style={ellipse, draw=accent, thick, fill=accent!10,
                minimum width=2.8cm, minimum height=1cm, text centered,
                text width=2.6cm},
  arrow/.style={-{Stealth[length=6pt]}, thick, color=primary!70},
  dasharrow/.style={-{Stealth[length=6pt]}, thick, dashed, color=gray}
]

% Capa nube
\node[cloud] (cloud) {Servidor central\\(nube,opcional)};

% Capa fog
\node[box, below=2cm of cloud] (fog1) {Nodo Fog\\Sector A};
\node[box, right=4cm of fog1]  (fog2) {Nodo Fog\\Sector B};
\node[box, left=2cm of fog1]   (fog3) {Nodo Fog\\Sector C};

% Capa campo
\node[box, below=1.6cm of fog3, fill=accent!8, draw=accent]
      (d1) {Brigada\\Rescate 1};
\node[box, below=1.6cm of fog1, fill=accent!8, draw=accent]
      (d2) {Brigada\\Rescate 2};
\node[box, right=0.6cm of d2,   fill=accent!8, draw=accent]
      (d3) {Dron\\Entrega};
\node[box, below=1.6cm of fog2, fill=accent!8, draw=accent]
      (d4) {Centro de\\Suministros};

% Flechas fog → cloud
\draw[dasharrow] (fog1.north) -- (cloud.south);
\draw[dasharrow] (fog2.north) -- (cloud.south);
\draw[dasharrow] (fog3.north) -- (cloud.south);

% Flechas Mesh fog ↔ fog
\draw[arrow, Stealth-Stealth] (fog1.east) -- (fog2.west);
\draw[arrow, Stealth-Stealth] (fog1.west) -- (fog3.east);

% Flechas campo → fog
\draw[arrow] (d1.north) -- (fog3.south);
\draw[arrow] (d2.north) -- (fog1.south);
\draw[arrow] (d3.north) -- (fog1.south);
\draw[arrow] (d4.north) -- (fog2.south);

% Etiquetas de protocolo
\node[gray, font=\scriptsize, above=0.15cm of fog1] {Mesh / Wi-Fi Direct};
\node[gray, font=\scriptsize, right=0.05cm of fog2.north east, rotate=30]
      {\scriptsize HTTPS};

\end{tikzpicture}
\caption{Arquitectura de tres capas de ResQ-Net. Las líneas continuas
  representan comunicación Mesh directa; las discontinuas, sincronización
  diferida a internet cuando está disponible.}
\label{fig:arch}
\end{figure}

\subsection{Flujo de datos}

\begin{enumerate}[leftmargin=1.5em]
  \item Una brigada registra en su tablet una \textbf{necesidad detectada}
        (p.\,ej., 12 personas con heridas graves en coordenada X).
  \item El dato se firma con timestamp y ID de nodo, y se replica por
        \textbf{Gossip Protocol} a todos los nodos fog del sector.
  \item Cada nodo fog ejecuta el \textbf{motor de triaje} y recalcula las
        prioridades de asignación de recursos disponibles.
  \item La decisión de asignación se devuelve a los dispositivos de campo como
        una instrucción concreta: ``Envíe kit médico tipo A a coordenada X.''
  \item Si hay internet, los datos se sincronizan con el servidor central para
        análisis post-desastre.
\end{enumerate}

\subsection{Componentes del MVP (simulado en navegador)}

El MVP web replica el comportamiento del sistema usando un mapa interactivo
(Leaflet.js) con nodos simulados. El motor de triaje corre enteramente en el
cliente (\textbf{JavaScript puro}), sin dependencia de servidor externo, lo
que demuestra la viabilidad de ejecutarlo en un nodo fog de bajo costo.

% ╔══════════════════════════════════════════════════════════════════════════╗
% ║  4. HERRAMIENTAS Y TECNOLOGÍAS                                           ║
% ╚══════════════════════════════════════════════════════════════════════════╝
\section{Herramientas y Tecnologías Utilizadas}

\begin{longtable}{>{\bfseries}p{3cm} p{4cm} p{6cm}}
  \toprule
  Tecnología & Categoría & Justificación \\
  \midrule
  \endfirsthead
  \toprule
  Tecnología & Categoría & Justificación \\
  \midrule
  \endhead
  HTML5 / CSS3 & Frontend & Estándar universal; corre en cualquier navegador
                            sin instalación adicional \\[4pt]
  JavaScript (ES2022) & Lógica cliente & Permite ejecutar el motor de triaje
                                          directamente en el navegador;
                                          sin dependencia de servidor \\[4pt]
  Leaflet.js 1.9 & Cartografía & Librería open-source de mapas interactivos;
                                  ligera (42\,KB), sin costo de API \\[4pt]
  OpenStreetMap & Datos de mapas & Cartografía libre y actualizada; no
                                    requiere clave de pago \\[4pt]
  PeerJS / WebRTC & Comunicación P2P & Conexión directa entre navegadores
                                        para simular red Mesh \\[4pt]
  Git / GitHub & Control de versiones & Trazabilidad completa del código;
                                         revisión entre integrantes \\[4pt]
  Overleaf & Documentación & Edición colaborativa de \LaTeX{} en tiempo real \\[4pt]
  Python 3.12 & Scripts auxiliares & Generación de datos de prueba y
                                      scripts de validación \\
  \bottomrule
\end{longtable}

\subsection{Justificación de tecnologías clave}

\textbf{JavaScript puro (sin framework):} elegimos no usar React ni Vue para
el MVP porque el objetivo es demostrar que el motor de triaje puede ejecutarse
en hardware de baja gama sin un sistema de empaquetado complejo. Un Raspberry
Pi 4 con Chromium puede correr la aplicación sin modificaciones.

\textbf{Leaflet.js en lugar de Google Maps:} Google Maps requiere una API key
de pago y conexión a sus servidores. Leaflet con tiles de OpenStreetMap puede
cargarse desde caché local, lo cual es crítico en escenarios sin internet.

% ╔══════════════════════════════════════════════════════════════════════════╗
% ║  5. METODOLOGÍA DE DESARROLLO                                            ║
% ╚══════════════════════════════════════════════════════════════════════════╝
\section{Metodología de Desarrollo}

Aplicamos \textbf{Scrum adaptado} en sprints de una semana, con la siguiente
estructura:

\begin{itemize}[leftmargin=1.5em]
  \item \textbf{Product Backlog:} lista priorizada de funcionalidades, desde
        el mapa interactivo hasta el motor de triaje.
  \item \textbf{Sprint Planning:} cada lunes (adaptado a miércoles/jueves
        según el cronograma del taller) definimos el sprint goal y distribuimos
        las tareas.
  \item \textbf{Daily standup (asíncrono):} mensajes diarios en el canal de
        WhatsApp del grupo con el formato: ¿Qué hice ayer? ¿Qué haré hoy?
        ¿Hay algún bloqueo?
  \item \textbf{Sprint Review:} cada avance ante el docente (13 y 20 de marzo)
        funcionó como revisión formal del sprint.
  \item \textbf{Retrospectiva:} al final de cada semana identificamos un punto
        de mejora concreto para el siguiente sprint.
\end{itemize}

El tablero de tareas se gestionó en \textbf{GitHub Projects} con columnas
\textit{Backlog}, \textit{In Progress}, \textit{In Review} y \textit{Done}.

% ╔══════════════════════════════════════════════════════════════════════════╗
% ║  6. PLANIFICACIÓN DE LAS 3 SEMANAS                                       ║
% ╚══════════════════════════════════════════════════════════════════════════╝
\section{Planificación de las 3 Semanas de Desarrollo}

\begin{longtable}{>{\bfseries}p{1.6cm} p{5cm} p{3cm} p{2.2cm}}
  \toprule
  Semana & Tarea & Responsable & Estado \\
  \midrule
  \endfirsthead
  \toprule
  Semana & Tarea & Responsable & Estado \\
  \midrule
  \endhead
  1 & Definición de arquitectura y tecnologías & Est.\,1 \& 3 & \textcolor{green!60!black}{Completado} \\
  1 & Estructura base HTML/CSS/JS & Est.\,2 & \textcolor{green!60!black}{Completado} \\
  1 & Boceto del algoritmo de triaje (pseudocódigo) & Est.\,3 & \textcolor{green!60!black}{Completado} \\
  1 & Inicio del documento \LaTeX & Est.\,4 & \textcolor{green!60!black}{Completado} \\
  1 & Configuración de GitHub y ramas & Est.\,5 & \textcolor{green!60!black}{Completado} \\
  \midrule
  2 & Mapa interactivo con nodos simulados & Est.\,2 & \textcolor{green!60!black}{Completado} \\
  2 & Motor de triaje v1 (puntuación básica) & Est.\,3 & \textcolor{green!60!black}{Completado} \\
  2 & Integración mapa + motor de triaje & Est.\,1 \& 2 & \textcolor{green!60!black}{Completado} \\
  2 & Secciones 1--4 del documento & Est.\,4 & \textcolor{green!60!black}{Completado} \\
  2 & Pruebas funcionales básicas & Est.\,5 & \textcolor{green!60!black}{Completado} \\
  \midrule
  3 & Motor de triaje v2 (múltiples recursos) & Est.\,3 & \textcolor{green!60!black}{Completado} \\
  3 & Panel de control y estadísticas & Est.\,2 & \textcolor{green!60!black}{Completado} \\
  3 & Secciones 5--10 del documento & Est.\,4 & \textcolor{green!60!black}{Completado} \\
  3 & Revisión integral y pruebas de estrés & Est.\,1 \& 5 & \textcolor{green!60!black}{Completado} \\
  3 & Preparación de la exposición & Todos & \textcolor{green!60!black}{Completado} \\
  \bottomrule
\end{longtable}

\subsection{Diagrama de Gantt simplificado}

\begin{figure}[H]
\centering
\begin{ganttchart}[
  hgrid,
  vgrid={*{4}{draw=none}, dotted},
  title/.style={fill=primary!80, draw=none, text=white},
  title label font=\small\bfseries,
  bar/.style={fill=primary!50, draw=primary},
  bar label font=\small,
  bar height=0.5,
  group/.style={fill=primary!20, draw=primary!60},
  milestone/.style={fill=accent, draw=accent!80},
  x unit=0.9cm,
  y unit title=0.6cm,
  y unit chart=0.55cm
]{1}{15}
  \gantttitle{Semana 1}{5}
  \gantttitle{Semana 2}{5}
  \gantttitle{Semana 3}{5} \\
  \gantttitlelist{1,...,15}{1} \\

  \ganttbar{Arquitectura}{1}{2} \\
  \ganttbar{HTML/CSS base}{1}{3} \\
  \ganttbar{Algoritmo triaje v1}{2}{5} \\
  \ganttbar{Documento \LaTeX{} (inicio)}{3}{5} \\
  \ganttmilestone{Avance 1}{5} \\
  \ganttbar{Mapa + nodos}{6}{8} \\
  \ganttbar{Integración}{7}{9} \\
  \ganttbar{Documento (secs.\,1--4)}{6}{10} \\
  \ganttmilestone{Avance 2}{10} \\
  \ganttbar{Triaje v2}{11}{12} \\
  \ganttbar{Panel de control}{11}{13} \\
  \ganttbar{Documento (secs.\,5--10)}{11}{14} \\
  \ganttbar{Pruebas finales}{13}{14} \\
  \ganttmilestone{Entrega final}{15}
\end{ganttchart}
\caption{Planificación Gantt del proyecto ResQ-Net.}
\label{fig:gantt}
\end{figure}

% ╔══════════════════════════════════════════════════════════════════════════╗
% ║  7. COMPONENTE ÉTICO                                                     ║
% ╚══════════════════════════════════════════════════════════════════════════╝
\section{Componente Ético}

\subsection{El dilema del triaje algorítmico}

ResQ-Net debe resolver un problema que ningún sistema técnico puede eludir:
\textit{¿a quién enviar el último kit de medicamentos cuando hay más personas
que recursos?} Esta pregunta tiene dimensiones étnicas, socioeconómicas y
médicas que un algoritmo simplifica necesariamente.

La tentación de diseñar un sistema que ``maximice vidas salvadas'' esconde
supuestos filosóficos con consecuencias concretas: ¿se priorizan los menores
de edad? ¿Se descarta a personas mayores de cierta edad? ¿Cómo se maneja la
incertidumbre en el diagnóstico de campo?

\textbf{Nuestra posición:} el algoritmo de triaje de ResQ-Net no toma
decisiones autónomas finales. Produce \textbf{recomendaciones priorizadas con
criterios explícitos y auditables} que siempre deben ser confirmadas por un
coordinador humano. La última palabra es siempre humana.

\subsection{Aplicación del Código de Ética IEEE/ACM}

Aplicamos principalmente el \textbf{Principio 1 del ACM Code of Ethics}:
\textit{``Un ingeniero de software actuará de manera consistente con el interés
público''} \cite{acm2018}, concretado en las siguientes decisiones de diseño:

\begin{itemize}[leftmargin=1.5em]
  \item \textbf{Transparencia algorítmica:} cada recomendación muestra los
        criterios usados (gravedad médica, distancia, stock disponible) para
        que el coordinador pueda cuestionarla.
  \item \textbf{Privacidad por defecto:} los datos de víctimas no contienen
        nombres propios ni identificadores personales; solo coordenadas,
        gravedad y tipo de necesidad.
  \item \textbf{Equidad:} el algoritmo no incluye variables socioeconómicas,
        étnicas ni de género. Solo criterios médicos y logísticos.
  \item \textbf{Trazabilidad:} cada decisión de asignación queda registrada
        con timestamp para auditoría post-crisis.
\end{itemize}

También aplicamos el \textbf{Principio 3} (\textit{Producto}): el MVP debe
ser seguro incluso en condiciones extremas. Por eso priorizamos la resiliencia
sobre la funcionalidad avanzada: es preferible un sistema que funcione sin
internet pero que pierda algunas características, a uno que falle por completo
cuando la conectividad no está disponible.

\subsection{Limitaciones éticas reconocidas}

\begin{itemize}[leftmargin=1.5em]
  \item El sistema depende de la calidad de los datos ingresados en campo
        (``garbage in, garbage out''). La fatiga y el estrés de las brigadas
        pueden introducir errores.
  \item La simulación del MVP no prueba el sistema bajo condiciones reales de
        caos, latencia de red o falla de hardware.
  \item El sesgo en los datos de entrenamiento (si se usa ML en versiones
        futuras) podría reproducir desigualdades históricas en la respuesta
        humanitaria.
\end{itemize}

% ╔══════════════════════════════════════════════════════════════════════════╗
% ║  8. IMPACTO SOCIAL ESPERADO                                              ║
% ╚══════════════════════════════════════════════════════════════════════════╝
\section{Impacto Social Esperado}

\subsection{Beneficiarios}

\begin{itemize}[leftmargin=1.5em]
  \item \textbf{Directos:} personas afectadas por desastres que reciben
        atención más rápida y coordinada; brigadas de rescate que operan con
        mejor información; coordinadores que toman decisiones con datos en
        tiempo real.
  \item \textbf{Indirectos:} ONGs y organismos estatales que analizan los
        datos post-crisis para mejorar protocolos futuros; aseguradoras y
        gobiernos que reducen costos de operaciones mal coordinadas.
\end{itemize}

\subsection{Indicadores de impacto}

\begin{longtable}{p{5.5cm} p{3cm} p{4cm}}
  \toprule
  Indicador & Línea base & Meta (6 meses) \\
  \midrule
  \endfirsthead
  \toprule
  Indicador & Línea base & Meta (6 meses) \\
  \midrule
  \endhead
  Tiempo promedio de primera respuesta coordinada & 4--6\,h & $<$\,1\,h \\
  Porcentaje de duplicación de entregas & 30--40\,\% & $<$\,5\,\% \\
  Cobertura de zonas sin asistencia & Sin métrica & $<$\,10\,\% sin visita \\
  Decisiones de triaje auditadas & 0\,\% & 100\,\% \\
  \bottomrule
\end{longtable}

\subsection{Escalabilidad}

La arquitectura Mesh permite incorporar nodos nuevos de forma orgánica:
basta con conectar un dispositivo a la red local del sector para que empiece
a sincronizar datos. No hay un servidor central que pueda convertirse en
cuello de botella. En una versión de producción, los nodos fog podrían ser
mochilas con Raspberry Pi y antenas de largo alcance (LoRa o Wi-Fi 5\,GHz),
desplegadas desde helicóptero en los primeros minutos del desastre.

% ╔══════════════════════════════════════════════════════════════════════════╗
% ║  9. CONCLUSIONES Y TRABAJO FUTURO                                        ║
% ╚══════════════════════════════════════════════════════════════════════════╝
\section{Conclusiones y Trabajo Futuro}

\subsection{Lecciones aprendidas}

Desarrollar ResQ-Net en tres semanas dejó lecciones que difícilmente habríamos
aprendido en el aula. La más importante: la arquitectura importa más que las
features. Dedicar la primera semana a diseñar la topología de nodos y el
protocolo de sincronización fue la decisión que salvó el proyecto; cada
funcionalidad que construimos después encajó sin fricciones estructurales.

Otra lección fue la dificultad de diseñar algoritmos que tomen decisiones
con consecuencias humanas reales. El debate interno sobre qué variables incluir
en el modelo de triaje ---y cuáles excluir para evitar discriminación--- fue
más útil que cualquier lectura sobre ética de IA.

\subsection{Limitaciones del MVP}

\begin{itemize}[leftmargin=1.5em]
  \item La red Mesh está simulada en el navegador; no probamos latencias ni
        pérdidas de paquetes reales.
  \item El mapa usa datos estáticos; en producción necesitaría actualizarse
        con imágenes satelitales post-desastre.
  \item El motor de triaje es determinista; en escenarios reales, la
        incertidumbre médica requeriría modelos probabilísticos.
\end{itemize}

\subsection{Trabajo futuro}

\begin{enumerate}[leftmargin=1.5em]
  \item \textbf{v2.0:} implementar la red Mesh real sobre dispositivos Android
        con Wi-Fi Direct y Bluetooth LE.
  \item \textbf{v2.1:} integrar imágenes de drones en tiempo real para
        actualizar el mapa automáticamente.
  \item \textbf{v3.0:} añadir un modelo de ML para predicción de demanda de
        recursos basado en el tipo y magnitud del desastre.
  \item \textbf{v3.1:} integración con sistemas de la UNGRD y la Cruz Roja
        Colombiana para pruebas piloto en simulacros reales.
\end{enumerate}

% ╔══════════════════════════════════════════════════════════════════════════╗
% ║  10. REFERENCIAS                                                         ║
% ╚══════════════════════════════════════════════════════════════════════════╝
\section*{Referencias}
\addcontentsline{toc}{section}{Referencias}

\begin{thebibliography}{99}

\bibitem{ocha2023}
United Nations Office for the Coordination of Humanitarian Affairs (OCHA).
\textit{Global Humanitarian Overview 2023}.
Ginebra: OCHA, 2023.
\href{https://www.unocha.org/global-humanitarian-overview-2023}{unocha.org/global-humanitarian-overview-2023}

\bibitem{ungrd2022}
Unidad Nacional para la Gestión del Riesgo de Desastres (UNGRD).
\textit{Informe anual de gestión del riesgo 2022}.
Bogotá: UNGRD, 2022.
\href{https://portal.gestiondelriesgo.gov.co}{portal.gestiondelriesgo.gov.co}

\bibitem{who2022}
World Health Organization (WHO).
\textit{Emergency Response Framework}, 2.ª ed.
Ginebra: WHO, 2022.
\href{https://www.who.int/publications/i/item/9789240058460}{who.int/publications}

\bibitem{schuller2016}
M. Schuller y P. Morales (eds.).
\textit{Tectonic Shifts: Haiti Since the Earthquake}.
Sterling: Kumarian Press, 2012.

\bibitem{plattner2010}
H. Plattner.
\textit{An Introduction to Design Thinking: Process Guide}.
Stanford: Hasso Plattner Institute of Design at Stanford University, 2010.

\bibitem{acm2018}
ACM \& IEEE-CS Joint Task Force on Software Engineering Ethics.
\textit{Software Engineering Code of Ethics and Professional Practice}, v5.2.
Nueva York: ACM, 2018.
\href{https://www.acm.org/code-of-ethics}{acm.org/code-of-ethics}

\bibitem{tanenbaum2017}
A. S. Tanenbaum y D. J. Wetherall.
\textit{Computer Networks}, 5.ª ed.
Upper Saddle River: Pearson, 2011.

\bibitem{bonomi2012}
F. Bonomi, R. Milito, J. Zhu y S. Addepalli.
``Fog Computing and its Role in the Internet of Things''.
En: \textit{Proc. of the First Edition of the MCC Workshop on Mobile Cloud Computing}, 2012, pp.\,13--16.

\bibitem{lee2020}
I. Lee y K. Lee.
``The Internet of Things (IoT): Applications, investments, and challenges for enterprises''.
\textit{Business Horizons}, vol.\,58, n.º\,4, pp.\,431--440, 2015.

\bibitem{openstreetmap}
OpenStreetMap contributors.
\textit{OpenStreetMap} [en línea].
Disponible en: \href{https://www.openstreetmap.org}{openstreetmap.org}
[Acceso: marzo 2026].

\end{thebibliography}

\end{document}